\documentclass[11pt,letterpaper]{report}

\usepackage[utf8]{inputenc}
\usepackage[T1]{fontenc}
\usepackage[spanish,provide=*]{babel}
\usepackage{amsmath,amssymb}
\usepackage{booktabs}
\usepackage{array}
\usepackage{graphicx}
\usepackage{float}
\usepackage[margin=2.5cm]{geometry}
\usepackage{enumitem}

\renewcommand{\arraystretch}{1.15}

\begin{document}

\begin{center}
{\LARGE\bfseries\linespread{1.3}\selectfont 1. INTRODUCCIÓN\par}
\end{center}
\vspace{0.5cm}

\setcounter{chapter}{1}
\setcounter{page}{11}

%==========================================================
Con el advenimiento de la computadora, una de las más importantes
herramientas para analizar el diseño y operación de sistemas o procesos
complejos es la simulación.

Aunque la construcción de modelos arranca desde el Renacimiento, el uso
moderno de la palabra simulación data de 1940, cuando los científicos
Von Neuman y Ulam que trabajaban en el proyecto Monte Carlo, durante la
Segunda Guerra Mundial, resolvieron problemas de reacciones nucleares
cuya solución experimental sería muy cara y el análisis matemático
demasiado complicado.

Con la utilización de la computadora en los experimentos de simulación,
surgieron incontables aplicaciones y con ello, una cantidad mayor de
problemas teóricos y prácticos. En este libro se intenta por
consiguiente, investigar y analizar cierto número de aplicaciones
importantes de simulación de las áreas de economía, administración de
negocios e investigación de operaciones, así como también sugerir
algunos métodos alternativos para resolver algunos problemas teóricos y
prácticos que surgen al efectuar simulaciones reales.

%==========================================================
\section*{1.1. Definición de Simulación}

Se ha empezado a utilizar la palabra simulación sin haber dado
previamente una definición de ella. Por consiguiente, antes de proseguir
con la discusión de este tema, sería conveniente describir algunas de
las definiciones más aceptadas y difundidas de la palabra simulación.
Thomas H. Naylor la define así:

\begin{quote}
\emph{Simulación es una técnica numérica para conducir experimentos en
una computadora digital. Estos experimentos comprenden ciertos tipos de
relaciones matemáticas y lógicas, las cuales son necesarias para
describir el comportamiento y la estructura de sistemas complejos del
mundo real a través de largos períodos de tiempo.}
\end{quote}

La definición anterior está en un sentido muy amplio, pues puede incluir
desde una maqueta, hasta un sofisticado programa de computadora. En
sentido más estricto, H. Maisel y G. Gnugnoli, definen simulación como:

\begin{quote}
\emph{Simulación es una técnica numérica para realizar experimentos en
una computadora digital. Estos experimentos involucran ciertos tipos de
modelos matemáticos y lógicos que describen el comportamiento de
sistemas de negocios, económicos, sociales, biológicos, físicos o
químicos a través de largos períodos de tiempo.}
\end{quote}

Otros estudiosos del tema como Robert E. Shannon, definen simulación
como:

\begin{quote}
\emph{Simulación es el proceso de diseñar y desarrollar un modelo
computarizado de un sistema o proceso y conducir experimentos con este
modelo con el propósito de entender el comportamiento del sistema o
evaluar varias estrategias con las cuales se puede operar el sistema.}
\end{quote}

Las definiciones anteriores no especifican si los sistemas modelados son
continuos o discretos. Sin embargo, es necesario señalar que el grueso
de este libro está dedicado al diseño, análisis y validación de sistemas
dinámicos discretos. Algunos autores de este tema como Geoffrey Gordon
en su libro \emph{System Simulation}, tratan a fondo el Análisis y
estudio de sistemas dinámicos continuos.

%==========================================================
\section*{1.2. Etapas para realizar un estudio de Simulación}

Se ha escrito mucho acerca de los pasos necesarios para realizar un
estudio de simulación. Sin embargo, la mayoría de los autores opinan que
los pasos necesarios para llevar a cabo un experimento de simulación
son:

\begin{itemize}[leftmargin=*]
\item \emph{Definición del sistema.} Para tener una definición exacta
del sistema que se desea simular, es necesario hacer primeramente un
análisis preliminar del mismo, con el fin de determinar la interacción
del sistema con otros sistemas, las restricciones del sistema, las
variables que interactúan dentro del sistema y sus interrelaciones, las
medidas de efectividad que se van a utilizar para definir y estudiar el
sistema y los resultados que se esperan obtener del estudio.

\item \emph{Formulación del modelo.} Una vez que están definidos con
exactitud los resultados que se esperan obtener del estudio, el
siguiente paso es definir y construir el modelo con el cual se
obtendrán los resultados deseados. En la formulación del modelo es
necesario definir todas las variables que forman parte de él, sus
relaciones lógicas y los diagramas de flujo que describan en forma
completa al modelo.

\item \emph{Colección de datos.} Es posible que la facilidad de
obtención de algunos datos o la dificultad de conseguir otros, pueda
influenciar el desarrollo y formulación del modelo. Por consiguiente, es
muy importante que se definan con claridad y exactitud los datos que el
modelo va a requerir para producir los resultados deseados. Normalmente,
la información requerida por un modelo se puede obtener de registros
contables, de órdenes de trabajo, de órdenes de compra, de opiniones de
expertos y si no hay otro remedio por experimentación.

\item \emph{Implementación del modelo en la computadora.} Con el modelo
definido, el siguiente paso es decidir si se utiliza algún lenguaje como
fortran, basic, algol, etc., o se utiliza algún paquete como GPSS,
simula, simscript, etc., para procesarlo en la computadora y obtener los
resultados deseados.

\item \emph{Validación.} Una de las principales etapas de un estudio de
simulación es la validación. A través de esta etapa es posible detallar
deficiencias en la formulación del modelo o en los datos alimentados al
modelo. Las formas más comunes de validar un modelo son:

\begin{enumerate}
\item La opinión de expertos sobre los resultados de la simulación.
\item La exactitud con que se predicen datos históricos.
\item La exactitud en la predicción del futuro.
\item La comprobación de falla del modelo de simulación al utilizar
datos que hacen fallar al sistema real.
\item La aceptación y confianza en el modelo de la persona que hará
uso de los resultados que arroje el experimento de simulación.
\end{enumerate}

\item \emph{Experimentación.} La experimentación con el modelo se
realiza después de que éste ha sido validado. La experimentación
consiste en generar los datos deseados y en realizar análisis de
sensibilidad de los índices requeridos.

\item \emph{Interpretación.} En esta etapa del estudio, se interpretan
los resultados que arroja la simulación y en base a esto se toma una
decisión. Es obvio que los resultados que se obtienen de un estudio de
simulación ayudan a soportar decisiones del tipo semi-estructurado, es
decir, la computadora en sí no toma la decisión, sino que la información
que proporciona ayuda a tomar mejores decisiones y por consiguiente a
sistemáticamente obtener mejores resultados.

\item \emph{Documentación.} Dos tipos de documentación son requeridos
para hacer un mejor uso del modelo de simulación. La primera se refiere
a la documentación de tipo técnico, es decir, a la documentación que el
departamento de Procesamiento de Datos debe tener del modelo. La segunda
se refiere al manual del usuario, con el cual se facilita la interacción
y el uso del modelo desarrollado, a través de una terminal de
computadora.
\end{itemize}

%==========================================================
\section*{1.3. Factores a considerar en el desarrollo del modelo de Simulación}

Puesto que la simulación está basada fuertemente en la teoría de
probabilidad y estadística, en matemáticas, en ciencias computacionales,
etc., es conveniente decir algunas ideas de cómo intervienen estas áreas
en el desarrollo y formulación del modelo de simulación.

\subsection*{1.3.1. Generación de variables aleatorias no-uniformes}

Si el modelo de simulación es estocástico, la simulación debe ser capaz
de generar variables aleatorias no-uniformes de distribuciones de
probabilidad teóricas o empíricas. Lo anterior puede ser obtenido si se
cuenta con un generador de números uniformes y una función que
transforme estos números en valores de la distribución de probabilidad
deseada. A este respecto, se han desarrollado una gran cantidad de
generadores para las distribuciones de probabilidad más comunes como: La
distribución normal, la distribución exponencial, la distribución
poisson, la distribución erlang, la distribución binomial, la
distribución gamma, la distribución beta, la distribución F, la
distribución t, etc.

\subsection*{1.3.2. Lenguajes de programación}

Las primeras etapas de un estudio de simulación se refieren a la
definición del sistema a ser modelado y a la descripción del sistema en
términos de relaciones lógicas de sus variables y diagramas de flujo.
Sin embargo, llega el momento de describir el modelo en un lenguaje que
sea aceptado por la computadora que se va a usar. En esta etapa se
tienen dos cursos de acción a seguir si no se tiene nada de software
sobre simulación: 1) Desarrollar el software requerido para estudios de
simulación, ó 2) Comprar software (lenguajes de programación de
propósito especial). Para esta alternativa es necesario analizar y
evaluar varios paquetes de simulación (GPSS, GASP, etc.) antes de tomar
la decisión final.

\subsection*{1.3.3. Condiciones iniciales}

La mayoría de los modelos de simulación estocástica se corren con la
idea de estudiar al sistema en una situación de estado estable. Sin
embargo, la mayoría de estos modelos presentan en su etapa inicial
estados transientes los cuales no son típicos del estado estable. Por
consiguiente es necesario establecer claramente las alternativas o
cursos de acción que existen para resolver este problema. Algunos
autores piensan que la forma de atacar este problema sería a través de:

\begin{itemize}[leftmargin=*]
\item Usar un tiempo de corrida lo suficientemente grande de modo que
los períodos transientes sean relativamente insignificantes con
respecto a la condición de estado estable.
\item Excluir una parte apropiada de la parte inicial de la corrida.
\item Utilizar simulación regenerativa.
\end{itemize}

Obviamente, de las tres alternativas presentadas, la que presenta menos
desventajas es el uso de simulación regenerativa. Las otras alternativas
presentan las desventajas de ser prohibitivamente excesivas en costo.

\subsection*{1.3.4. Tamaño de la muestra}

Uno de los factores principales a considerar en un estudio de simulación
es el tamaño de la muestra (número de corridas en la computadora). La
selección de un tamaño de muestra apropiado que asegure un nivel deseado
de precisión y a la vez minimice el costo de operación del modelo, es un
problema algo difícil pero muy importante. Puesto que la información
proporcionada por el experimento de simulación sería la base para
decidir con respecto a la operación del sistema real, esta información
deberá ser tan exacta y precisa como sea posible o al menos el grado de
imprecisión presente en la información proporcionada por el modelo debe
ser conocida. Por consiguiente, es necesario que un análisis estadístico
sea realizado para determinar el tamaño de muestra requerido.

El tamaño de la muestra puede ser obtenido de dos maneras:

\begin{enumerate}
\item Previa e independientemente de la operación del modelo, o
\item Durante la operación del modelo y basado en los resultados
arrojados por el modelo. Para la última alternativa se utiliza la
técnica estadística de intervalos de confianza.
\end{enumerate}

\subsection*{1.3.5. Diseño de experimentos}

El diseño de experimentos es un tópico cuya relevancia en experimentos
de simulación ha sido reconocida pero raramente aplicada. El diseño de
experimentos en estudios de simulación puede ser de varios tipos,
dependiendo de los propósitos específicos que se hayan planteado.
Existen varios tipos de análisis que pueden ser requeridos. Entre los
más comunes e importantes se pueden mencionar los siguientes:

\begin{itemize}[leftmargin=*]
\item Comparación de las medias y varianzas de las alternativas
analizadas.
\item Determinación de la importancia y el efecto de diferentes
variables en los resultados de la simulación.
\item Búsqueda de los valores óptimos de un conjunto de variables.
\end{itemize}

Para realizar el primer tipo de análisis, al cual se le denomina
comúnmente diseño de experimentos de un factor simple, es necesario
tomar muy en cuenta el tamaño de la muestra, las condiciones iniciales y
la presencia o ausencia de autocorrelación. Para el segundo tipo de
análisis, existe una gran cantidad de literatura, puesto que la gran
mayoría de los libros de texto de diseño de experimentos, explican o
tratan el tema de análisis de varianza y técnicas de regresión como
medios para evaluar la importancia y el efecto de varias variables en
los resultados de operación de un sistema. Para el tercer tipo de
análisis, generalmente se requiere utilizar algoritmos heurísticos de
búsqueda como por ejemplo el algoritmo de Hooke y Jeeves.

%==========================================================
\section*{1.4. Ventajas y desventajas en el uso de Simulación}

Aunque la técnica de simulación generalmente se ve como un método de
último recurso, recientes avances en las metodologías de simulación y la
gran disponibilidad de software que actualmente existe en el mercado,
han hecho que la técnica de simulación sea una de las herramientas más
ampliamente usadas en el análisis de sistemas. Además de las razones
antes mencionadas, Thomas H. Naylor ha sugerido que un estudio de
simulación es muy recomendable porque presenta las siguientes ventajas:

\begin{itemize}[leftmargin=*]
\item A través de un estudio de simulación, se puede estudiar el efecto
de cambios internos y externos del sistema, al hacer alteraciones en el
modelo del sistema y observando los efectos de esas alteraciones en el
comportamiento del sistema.
\item Una observación detallada del sistema que se está simulando puede
conducir a un mejor entendimiento del sistema y por consiguiente a
sugerir estrategias que mejoren la operación y eficiencia del sistema.
\item La técnica de simulación puede ser utilizada como un instrumento
pedagógico para enseñar a estudiantes habilidades básicas en análisis
estadístico, análisis teórico, etc.
\item La simulación de sistemas complejos puede ayudar a entender mejor
la operación del sistema, a detectar las variables más importantes que
interactúan en el sistema y a entender mejor las interrelaciones entre
estas variables.
\item La técnica de simulación puede ser usada para experimentar con
nuevas situaciones, sobre las cuales se tiene poca o ninguna información.
A través de esta experimentación se puede anticipar mejor a posibles
resultados no previstos.
\item La técnica de simulación se puede utilizar también para
entrenamiento de personal. En algunas ocasiones se puede tener una buena
representación de un sistema (como por ejemplo los juegos de negocios),
y entonces a través de él es posible entrenar y dar experiencia a cierto
tipo de personal.
\item Cuando nuevos elementos son introducidos en un sistema, la
simulación puede ser usada para anticipar cuellos de botella o algún
otro problema que puede surgir en el comportamiento del sistema.
\end{itemize}

A diferencia de las ventajas mencionadas, la técnica de simulación
presenta el problema de requerir equipo computacional y recursos humanos
costosos. Además, generalmente se requiere bastante tiempo para que un
modelo de simulación sea desarrollado y perfeccionado. Finalmente, es
posible que la alta administración de una organización no entienda esta
técnica y esto crea dificultad en vender la idea.

%==========================================================
\section*{1.5. Ejemplos de usos de Simulación}

Existe una gran cantidad de áreas donde la técnica de simulación puede
ser aplicada. Algunos ejemplos podrían ser los siguientes:

\begin{itemize}[leftmargin=*]
\item \emph{Simulación de un sistema de colas.} Con la técnica de
simulación es posible estudiar y analizar sistemas de colas cuya
representación matemática sería demasiado complicada de analizar.
Ejemplos de estos sistemas serían aquellos donde es posible la llegada
al sistema en grupo, la salida de la cola del sistema, el rehusar entrar
al sistema cuando la cola es excesivamente grande, etc.
\item \emph{Simulación de un sistema de inventarios.} A través de
simulación se pueden analizar más fácilmente sistemas de inventarios
donde todos sus parámetros (tiempo de entrega, demanda, costo de llevar
inventario, etc.), son estocásticos.
\item \emph{Simulación de un proyecto de inversión.} Existen en la
práctica una gran cantidad de proyectos de inversión donde la
incertidumbre con respecto a los flujos de efectivo que el proyecto
genera a las tasas de interés, a las tasas de inflación, etc., hacen
difícil y a veces imposible manejar analíticamente este tipo de
problemas. Para este tipo de situaciones el uso de simulación es
ampliamente recomendado.
\item \emph{Simulación de sistemas económicos.} La técnica de
simulación puede ser utilizada para evaluar el efecto de cierto tipo de
decisiones (devaluación de la moneda, el impuesto al valor agregado,
etc.), en las demás variables macroeconómicas como: producto nacional
bruto, balanza comercial, inflación, oferta monetaria, circulante, etc.
\item \emph{Simulación de estados financieros.} La expansión y
diversificación de una organización a través de la adquisición y
establecimiento de nuevas empresas, repercuten significativamente en su
posición y estructura financiera. Por consiguiente, el uso de simulación
permite analizar cuál de las estrategias de crecimiento son las que
llevarán a la organización al logro de sus objetivos y metas de corto,
mediano y largo plazos.
\end{itemize}

\end{document}
